%!TEX root = ../main.tex
\section{Discussion and Lessons Learned}
\label{sec:discussion}

\subsection{What worked}

The most reliable modeling choice was \emph{constrained complementarity}.
G4 and CQT emphasize different physical views: energy/pausing versus
log-frequency spectral structure. Fixed equal score fusion adds no
parameter that a tiny calibration subset can overfit. Its advantage persisted
across six Train-only partitions and both whole-side transfer directions,
although CQT's directional variability warns against claiming universal fusion
dominance. Path signatures supplied a more orthogonal temporal view and
produced the best hidden-Test UAR. The deep global/local systems reached nearly
the same UAR while producing a less cold-aggressive operating point.

More important, separating model selection, threshold selection and final
evaluation worked. Whole-official-side transfer avoids depending on inferred
identities for the load-bearing comparison. Repeated partitions and paired
resampling revealed variation that model-seed standard deviations had hidden.
Reporting both class recalls or secondary metrics exposed behavior that UAR
alone concealed.

\subsection{What failed and why}

Our largest failure was treating a repeatedly consulted Development subset as
a holdout. Once architectures, features, fusion weights and thresholds had all
been compared there, its score measured cumulative selection as well as model
quality. Optimization-seed variation ($\approx0.006$ for the historical K=2
system) was also much smaller than boundary variation ($\approx0.017$), so it
was the wrong uncertainty axis.

The speaker-proxy audit produced a second, more general lesson: validating a
representation is not the same as validating labels transferred to a new
domain. Side-local ECAPA clusters were reproducible, but nearest-Train-centroid
assignments fragmented Development-local structure. A split can therefore
have zero overlap under the labels used to construct it while still failing an
independent held-side audit. Future proxy construction must be label-free,
side-local and tested for boundary overlap before it is used in a performance
claim.

Finally, flexibility was usually a liability at this sample size. Learned
stacking, complex foundation-model heads, augmentation and explicit
de-confounding supplied more ways to fit a small number of cold participants
or properties of the preprocessing operation. Negative controls were more
informative than further hyperparameter search, especially self-splicing,
matched bottlenecks and discriminator train/held gaps.

\subsection{What we would do differently}

If restarting the project, we would write and freeze the evaluation plan before
training a large model. Architecture, feature, epoch and threshold selection
would occur only in nested grouped CV inside official Train. The speaker
embedding, clustering algorithm and $k$ would be chosen without cold labels;
the resulting group assignments would then be challenged side-locally through
method agreement, stability and explicit boundary-overlap tests. Official
Development would be used once after the full pipeline had been locked. Only
then would we refit on Train+Development and create Test predictions.

We would also distinguish three uncertainty questions from the beginning:
optimization variability across model seeds, selection variability across
folds, and population variability across held-out participants. The last is
the health-relevant quantity, but it cannot be recovered from seed standard
deviation. Every comparison would therefore use paired out-of-fold predictions
and participant-level resampling where verified identities are available. The
decision threshold and the relative cost of missed colds versus false alarms
would be declared from the intended use case rather than selected solely for
the leaderboard metric.

\subsection{Implications for health ML}

The effective sample is the participant, not the audio chunk. Thirty-seven
cold participants per official fitting partition cannot support fine-grained
claims from thousandths of UAR. Moreover, high UAR is not a clinical operating
policy: false negatives and false positives have different consequences, and
Submissions 3 and 5 demonstrate that the same UAR can mask very different
decision behavior. A future study should predefine the target use case,
participant-level uncertainty, threshold policy and external validation before
optimizing architecture.

\subsection{Limitations}

True URTIC speaker IDs are absent from the student release; all identity
analyses therefore use imperfect proxies. Cold status is a thresholded
self-report and each participant is recorded in only one state, so identity,
session and health effects cannot be causally separated. Development received
substantial historical selection exposure, and the five hidden-Test results
are single leaderboard observations. We did not perform prospective clinical
calibration, longitudinal within-person evaluation or external-corpus
validation. Our negative results apply only to the tested feature extractors,
heads and de-confounding recipes; they are not evidence that the broader model
families cannot work.

\subsection{Reproducibility, ethics and disclosure}

The project retains feature caches, model configurations, prediction files and
machine-readable result summaries for every number reported here. In
particular, the evaluation audit, held-side overlap audit, repeated G4+CQT CV,
whole-side transfer and threshold-transfer analyses are stored as separate
artifacts rather than reconstructed from presentation slides. Submission 5's
use of unlabeled Test features for CORAL is disclosed because it changes the
evaluation setting even though it uses no labels.

Voice is both health data and a biometric signal. A system that unintentionally
recognizes speakers can expose identity while appearing to predict health.
The present results therefore support neither deployment nor individual
diagnosis. Any follow-up should use consent appropriate to biometric and health
processing, examine performance across demographic and recording-condition
subgroups, minimize retention of identifiable audio, and validate prospectively
against a clinically meaningful reference rather than only self-report.
